% !TeX program = xelatex
\documentclass{cjlu_proposal}
\usepackage[backend=biber,style=gb7714-2015,sorting=none]{biblatex}
\usepackage[colorlinks=true,linkcolor=black,citecolor=blue,urlcolor=blue]{hyperref}
\addbibresource{review_refs.bib}

\studentname{李康峰}
\studentid{2201400216}
\major{自动化}
\classinfo{22 工试\hspace{0.5em}2 班}
\thesistitle{基于人工智能的缺陷射线图像识别与检测系统设计}
\advisor{洪宇翔}
\school{机电工程学院}
\coverdate{2026 年 1 月 12 日}

\begin{document}

\makeopeningcover

\section{选题的背景与意义}
\subsection{你还}
说明选题的动机、与学科前沿的关系以及实际/工程价值；可引用综述来支撑必要性\cite{zhang2024review}。

\section{国内外发展现状}
列出已有研究成果、技术路线及局限，展示尚未解决的问题\cite{wang2023method}。

\section{研究的基本内容与拟解决的主要问题}
\begin{enumerate}[leftmargin=*]
  \item 重点研究内容一：描述具体模型/系统/子模块。
  \item 重点研究内容二：说明预期达到的指标或性能。
  \item 拟解决的关键问题：明确难点（例如数据稀缺、实时性、可解释性等）。
\end{enumerate}

\section{技术路线、研究方案与可行性分析}
\begin{enumerate}[leftmargin=*]
  \item 系统整体设计：整体结构图+模块划分。
  \item 硬件选型与部署：说明所用传感器/控制器等。
  \item 软件流程：数据处理、算法推理、接口流程。
  \item 可行性分析：分别分析硬件/软件可行性。
\end{enumerate}

\section{实施计划}
\begin{tabularx}{\textwidth}{p{6.5em}X}
  \toprule
  周次        & 内容                         \\
  \midrule
  第 12-13 周 & 确定课题方向，整理需求背景，制定实验计划       \\
  第 13-14 周 & 明确关键指标与评价方法，完成任务书与预研文档     \\
  第 15-19 周 & 查阅文献、搭建实验平台、撰写初稿及中期报告      \\
  第 1-2 周   & 在导师指导下对可行性分析与初步成果修订，准备答辩材料 \\
  第 3-5 周   & 实施详细硬件电路与控制逻辑、进行算法程序的实现    \\
  第 6 周     & 完成毕业设计中期检查、整理阶段性成果与问题清单    \\
  第 7-9 周   & 深入调优算法、完成整体调试与性能验证         \\
  第 10 周    & 参与毕业设计验收、提交最终报告            \\
  第 11-12 周 & 撰写论文、导师审核、定稿打印             \\
  第 13 周    & 接受评阅、总结经验、准备答辩             \\
  第 14-15 周 & 完成毕业答辩、归档设计材料              \\
  \bottomrule
\end{tabularx}

\section{参考文献}
文末列表按照学校要求格式，引用条目不少于20篇，英文文献不少于5篇（含需翻译的两篇）。下表摘自模板，供整理文献类型时参考。

\begin{tabularx}{\textwidth}{p{0.18\textwidth}X}
  \toprule
  文献类型           & 格式与示例                                                                                                                                                           \\
  \midrule
  专著、论文集、学位论文、报告 & [序号]主要责任者．文献题名[文献类型标识]．版次（第1版可省略）．出版地:出版者,出版年．起止页码． 示例：[1]刘国钧等．图书馆目录[M]．北京:高等教育出版社,1957．15-18．                                                                  \\
  期刊文章           & [序号]主要责任者．文献题名[文献类型标识]．刊名,年,卷(期):起止页码． 示例：[2]金显贺,王昌长,王忠东等．一种用于在线检测局部放电的数字滤波技术[J]．清华大学学报(自然科学版),1993,33(4):62-67．                                                \\
  析出文献           & [序号]析出文献主要责任者．析出文献题名[文献类型标识]．原文献主要责任者．原文献题名[C]．出版地:出版者,出版年．析出文献起止页码． 示例：[3]钟文发．非线性规划在可燃毒物配置中的应用[A]．赵玮．运筹学的理论与应用——中国运筹学会第五届大会论文集[C]．西安:西安电子科技大学出版社,1996．468-471． \\
  国际、国家标准名称      & [序号]标准编号,标准名称[文献类型标识]． 示例：[4]GB/T 16159-1996,汉语拼音正词法基本规则[S]．                                                                                                    \\
  专利             & [序号]专利所有者．专利题名[文献类型标识]．专利国别:专利号,出版日期． 示例：[5]姜锡洲.一种温热外敷药制备方案[P]. 中国专利:881056073,1989-07-26．                                                                      \\
  报纸文章           & [序号]主要责任者．文献题名[文献类型标识]．报纸名,出版日期（版次）． 示例：[6]谢希德.创造学习的新思路[N]. 人民日报,1998-12-25(10).                                                                                \\
  电子文献           & [序号]主要责任者．文献题名[文献类型标识]．电子文献出处或可获取地址,发表/更新或引用日期． 示例：[7]王明亮．关于中国学术期刊标准化数据库系统工程的进展[EB/OL]．http://www.cajcd.edu.cn/pub/wml.txt/980810-2.html,1998-08-16/1998-10-04． \\
  \bottomrule
\end{tabularx}

\printbibliography[heading=none]

\end{document}
